\chapter{Noise and Error}
\lecture{8}{30 Sep.}{}

\begin{prev}
    \autoref{sec:vc-interpret} ended with the VC message: learning happens if
    \(\dvc\) is finite, \(N\) is large, and \(\Ein\) is low. The whole
    derivation, however, assumed a deterministic target \(f\), labels
    \(y_n = f(\bm{x}_n)\) with no noise, and binary classification measured by
    the fraction of mistakes.
\end{prev}

This chapter relaxes the first two assumptions and questions the third. Noise
turns the target function into a target \emph{distribution}, and the VC bound
survives the change. Once the target is a distribution, `the best prediction'
depends on how errors are measured, so the error measure becomes a component of
the learning problem in its own right --- chosen by the application, and then
often replaced by something the algorithm can actually optimise. The last
section works one such case, weighted classification, end to end.

\section{Noise and Probabilistic Target}\label{sec:noise}

\source{slides p.2--7}

\subsection{Recap: The Learning Flow}

\begin{center}
    \begin{adjustbox}{max width=\textwidth}
        \begin{tikzpicture}
            \node[mlbox, minimum width=4.4cm] (f) at (0,2.5)
                {unknown target function\\
                 \(f\colon \mathcal{X} \to \mathcal{Y}\)\\
                 \colorbox{WildStrawberry!20}{\(+\) noise}\\[2pt]
                 {\scriptsize (ideal credit approval formula)}};
            \node[mlbox, minimum width=2.6cm] (P) at (5.4,2.5)
                {unknown\\\(P\) on \(\mathcal{X}\)};
            \node[mlbox, minimum width=5.2cm] (D) at (0,0)
                {training examples\\
                 \(\mathcal{D}\colon (\bm{x}_1, y_1), \dots , (\bm{x}_N, y_N)\)\\[2pt]
                 {\scriptsize (historical records in bank)}};
            \node[mlop, minimum width=2.4cm] (A) at (5.4,0)
                {learning\\algorithm\\\(\mathcal{A}\)};
            \node[mlbox, minimum width=4.4cm] (g) at (10.2,0)
                {final hypothesis\\\(g \approx f\)\\[2pt]
                 {\scriptsize (`learned' formula to be used)}};
            \node[mlbox, minimum width=3.4cm] (H) at (5.4,-2.4)
                {hypothesis set\\\(\mathcal{H}\)\\[2pt]
                 {\scriptsize (set of candidate formulas)}};
            \draw[mlarrow] (f) -- (D);
            \draw[mlarrow, dashed] (P) -- (D)
                node[pos=0.58, font=\scriptsize, fill=white, inner sep=1.5pt]
                {\(\bm{x}_1, \bm{x}_2, \dots , \bm{x}_N\)};
            \draw[mlarrow, dashed] (P.east) -| (g.north)
                node[pos=0.25, above, font=\scriptsize] {\(\bm{x}\)};
            \draw[mlarrow] (D) -- (A);
            \draw[mlarrow] (A) -- (g);
            \draw[mlarrow] (H) -- (A);
        \end{tikzpicture}
    \end{adjustbox}
\end{center}

This is the flow of \autoref{sec:bin-learning}, with one addition in the target
box.

\begin{question}
    What if there is noise?
\end{question}

\subsection{Noise}

\begin{center}
    \begin{minipage}[c]{0.3\textwidth}
        \centering
        \begin{tikzpicture}[scale=0.7]
            \fill[poshalf] (0,3.2) -- (3.2,0) -- (3.2,3.2) -- cycle;
            \fill[neghalf] (0,3.2) -- (3.2,0) -- (0,0) -- cycle;
            \draw[sepline] (0,3.2) -- (3.2,0);
            \draw[black!45] (0,0) rectangle (3.2,3.2);
            \foreach \p in {(1.5,2.4), (2.2,2.8), (2.6,1.6), (1.9,1.9),
                (2.9,2.4), (0.4,3.0)}{\node[dpos] at \p {};}
            \foreach \p in {(0.5,0.9), (1.2,1.0), (0.8,1.9), (1.8,0.6),
                (0.4,2.0), (2.7,0.3)}{\node[dneg] at \p {\(\times\)};}
            \node[dneg] at (2.4,2.1) {\(\times\)};
            \node[dneg] at (1.3,2.9) {\(\times\)};
            \node[dpos] at (0.9,0.7) {};
            \node[dpos] at (0.3,0.25) {};
        \end{tikzpicture}
    \end{minipage}%
    \begin{minipage}[c]{0.62\textwidth}
        Noise was briefly introduced before the \textbf{pocket} algorithm
        (\autoref{sec:nonsep}), as the reason a data set might not be linearly
        separable even when the underlying rule is a line: a few examples sit
        on the wrong side of the best line.
    \end{minipage}
\end{center}

But there is more! Take the credit applicant again:

\begin{center}
    \begin{minipage}[c]{0.4\textwidth}
        \centering
        \small
        \begin{tabular}{@{}|c|c|@{}}
            \hline
            age & 23 years\\
            gender & female\\
            annual salary & NTD 1,000,000\\
            year in residence & 1 year\\
            year in job & 0.5 year\\
            current debt & 200,000\\
            \hline
        \end{tabular}\\[3pt]
        credit? \(\left\{ \text{no}(-1), \text{yes}(+1) \right\}\)
    \end{minipage}%
    \begin{minipage}[c]{0.55\textwidth}
        \begin{itemize}
            \item \textbf{noise in \(y\)}: good customer, `mislabeled' as bad?
            \item \textbf{noise in \(y\)}: same customers, different labels?
            \item \textbf{noise in \(\bm{x}\)}: inaccurate customer
                information?
        \end{itemize}
    \end{minipage}
\end{center}

The first is a clerical error on the label; the second is two applicants with
identical records but different credit histories, which no function of
\(\bm{x}\) can fit; the third corrupts the input rather than the label. In
every case the pair \((\bm{x}, y)\) is no longer tied by a function \(y =
f(\bm{x})\).

\begin{question}
    Does the VC bound work under noise?
\end{question}

\subsection{Probabilistic Marbles}

One key of the VC bound is the marbles of \autoref{sec:hoeffding}:

\begin{center}
    \begin{tikzpicture}
        \begin{scope}
            \clip[scale=1.3] \binpath -- cycle;
            \foreach \r in {0,...,13}{
                \foreach \c in {0,...,11}{
                    \pgfmathtruncatemacro{\m}{mod(\r*7 + \c*5 + \r*\c, 4)}
                    \pgfmathsetmacro{\xx}{-0.62 + 0.113*\c + 0.056*mod(\r,2)}
                    \pgfmathsetmacro{\yy}{0.05 + 0.093*\r}
                    \ifnum\m=0
                        \node[morange=3.2pt] at (\xx,\yy) {};
                    \else
                        \node[mgreen=3.2pt] at (\xx,\yy) {};
                    \fi}}
        \end{scope}
        \draw[binwall, scale=1.3] \binpath;
        \node[font=\scriptsize\bfseries, below] at (0,-0.02) {bin};
        \draw[scoop] (0.1,1.35) .. controls (0.8,2.0) and (1.8,1.8) ..
            (2.0,0.85);
        \node[font=\scriptsize\bfseries] at (1.0,2.0) {sample};
        \foreach \i/\c in {0/mgreen,1/mgreen,2/morange,3/mgreen,4/mgreen,
            5/mgreen,6/morange,7/mgreen,8/mgreen,9/morange}{
            \node[\c=3.4pt] at (1.55+0.14*\i,0.65) {};}
    \end{tikzpicture}
\end{center}

Noise changes how a marble gets its colour, and nothing else:

\begin{center}
    \solidrules
    \begin{tabular}{@{}l|l@{}}
        \toprule
        `deterministic' marbles & `probabilistic' (noisy) marbles\\
        \midrule
        marble \(\bm{x} \sim P(\bm{x})\) & marble \(\bm{x} \sim P(\bm{x})\)\\
        deterministic colour \(\llbracket f(\bm{x}) \neq h(\bm{x})
        \rrbracket\)
        & probabilistic colour \(\llbracket y \neq h(\bm{x}) \rrbracket\)
        with \(y \sim P(y \mid \bm{x})\)\\
        \bottomrule
    \end{tabular}
\end{center}

A noisy marble is one whose colour is decided afresh, by a coin flip, each time
it is drawn. Hoeffding's inequality never cared how the colour was produced:
all it needs is that the draws are i.i.d.\ and each is orange with the same
probability.

\begin{moral}
    Same nature: can estimate \(\mathbb{P}[\text{orange}]\) if
    \(\overset{\text{i.i.d.}}{\sim}\).
\end{moral}

So the whole chain of the last three chapters goes through if the examples are
drawn as
\[
    \text{VC holds for }
    \underbrace{\bm{x} \overset{\text{i.i.d.}}{\sim} P(\bm{x}), \;
    y \overset{\text{i.i.d.}}{\sim} P(y \mid \bm{x})}_{(\bm{x}, y)
    \overset{\text{i.i.d.}}{\sim} P(\bm{x}, y)} .
\]
The out-of-sample error of \(h\) now reads \(\Eout(h) = \mathbb{P}_{(\bm{x},
y)}\left[ h(\bm{x}) \neq y \right]\), the orange fraction of the noisy bin.

\subsection{Target Distribution \texorpdfstring{\(P(y \mid \bm{x})\)}{P(y|x)}}

\begin{definition}[Target distribution]\label{def:target-dist}
    The \vocab{target distribution} \(P(y \mid \bm{x})\) characterizes the
    behaviour of the `mini-target' on one \(\bm{x}\).
\end{definition}

\begin{itemize}
    \item It can be viewed as `ideal mini-target' \(+\) noise, e.g.
    \begin{itemize}
        \item \(P(\circ \mid \bm{x}) = 0.7\), \(P(\times \mid \bm{x}) = 0.3\);
        \item ideal mini-target \(f(\bm{x}) = \circ\);
        \item `flipping' noise level \(= 0.3\).
    \end{itemize}
    \item A deterministic target \(f\) is a \textbf{special case of target
        distribution}:
    \begin{itemize}
        \item \(P(y \mid \bm{x}) = 1\) for \(y = f(\bm{x})\);
        \item \(P(y \mid \bm{x}) = 0\) for \(y \neq f(\bm{x})\).
    \end{itemize}
\end{itemize}

So nothing is lost by switching: every earlier setting is a target distribution
concentrated on one label. The two distributions divide the work:

\begin{moral}
    Goal of learning: predict the ideal mini-target (w.r.t.\ \(P(y \mid
    \bm{x})\)) on often-seen inputs (w.r.t.\ \(P(\bm{x})\)).
\end{moral}

\(P(y \mid \bm{x})\) is what we try to learn; \(P(\bm{x})\) only decides where
it matters to get it right.

\subsection{The New Learning Flow}

In the flow, the target function box becomes a target distribution, which now
supplies the labels \(y_1, \dots , y_N\) of the training set \emph{and} the
label \(y\) against which \(g\) is judged on a future \(\bm{x}\).

\begin{center}
    \begin{adjustbox}{max width=\textwidth}
        \begin{tikzpicture}
            \node[mlbox, minimum width=4.4cm, draw=orange!85!black,
                fill=orange!8] (f) at (0,2.5)
                {unknown target\\
                 distribution \(P(y \mid \bm{x})\)\\
                 containing \(f(\bm{x})\) \(+\) noise\\[2pt]
                 {\scriptsize (ideal credit approval formula)}};
            \node[mlbox, minimum width=2.6cm] (P) at (5.4,2.5)
                {unknown\\\(P\) on \(\mathcal{X}\)};
            \node[mlbox, minimum width=5.2cm] (D) at (0,0)
                {training examples\\
                 \(\mathcal{D}\colon (\bm{x}_1, y_1), \dots , (\bm{x}_N, y_N)\)\\[2pt]
                 {\scriptsize (historical records in bank)}};
            \node[mlop, minimum width=2.4cm] (A) at (5.4,0)
                {learning\\algorithm\\\(\mathcal{A}\)};
            \node[mlbox, minimum width=4.4cm] (g) at (10.2,0)
                {final hypothesis\\\(g \approx f\)\\[2pt]
                 {\scriptsize (`learned' formula to be used)}};
            \node[mlbox, minimum width=3.4cm] (H) at (5.4,-2.4)
                {hypothesis set\\\(\mathcal{H}\)\\[2pt]
                 {\scriptsize (set of candidate formulas)}};
            \draw[mlarrow, dashed, draw=orange!85!black] (f) -- (D)
                node[midway, left, font=\scriptsize, text=orange!85!black]
                {\(y_1, y_2, \dots , y_N\)};
            \draw[mlarrow, dashed, draw=orange!85!black] (f.north)
                .. controls (1,4.4) and (11.1,4.4) .. ([xshift=0.9cm]g.north)
                node[pos=0.5, above, font=\scriptsize, text=orange!85!black]
                {\(y\)};
            \draw[mlarrow, dashed] (P) -- (D)
                node[pos=0.58, font=\scriptsize, fill=white, inner sep=1.5pt]
                {\(\bm{x}_1, \bm{x}_2, \dots , \bm{x}_N\)};
            \draw[mlarrow, dashed] (P.east) -| ([xshift=-0.9cm]g.north)
                node[pos=0.3, above, font=\scriptsize] {\(\bm{x}\)};
            \draw[mlarrow] (D) -- (A);
            \draw[mlarrow] (A) -- (g);
            \draw[mlarrow] (H) -- (A);
        \end{tikzpicture}
    \end{adjustbox}
\end{center}

\begin{moral}
    VC still works, pocket algorithm explained :-)
\end{moral}

Pocket was introduced for data that no line fits. With noise, that is the
normal case even when the ideal mini-target is linear; the VC bound says the
line with the smallest \(\Ein\) still has \(\Eout \approx \Ein\), so pocket's
choice is justified.

\begin{exercise}[Fun Time]
    Let's revisit PLA/pocket. Which of the following claim is true?
    \begin{enumerate}[(1)]
        \item In practice, we should try to compute if \(\mathcal{D}\) is
            linear separable before deciding to use PLA.
        \item If we know that \(\mathcal{D}\) is not linear separable, then the
            target function \(f\) must not be a linear function.
        \item If we know that \(\mathcal{D}\) is linear separable, then the
            target function \(f\) must be a linear function.
        \item None of the above.
    \end{enumerate}
\end{exercise}

\begin{proof}[Reference Answer]
    (4). (1) After computing if \(\mathcal{D}\) is linear separable, we shall
    know \(\bm{w}^*\) and then there is no need to use PLA. (2) What about
    noise? A linear target plus flipping noise gives non-separable data.
    (3) What about `sampling luck'? :-) A curved target can happen to produce a
    sample that some line separates.
\end{proof}

\section{Error Measure}\label{sec:error-measure}

\source{slides p.8--13}

The target is now a distribution, and \(g\) is supposed to be \(\approx f\).
Before learning anything, we need to say what `\(\approx\)' means.

\subsection{Error Measure}

\begin{center}
    \tikz\node[mlbox, minimum width=3cm] {final hypothesis\\\(g \approx f\)};
\end{center}

\begin{itemize}
    \item How well? Previously, we considered the out-of-sample measure
        \[
            \Eout(g) = \mathop{\mathcal{E}}_{\bm{x} \sim P}
            \llbracket g(\bm{x}) \neq f(\bm{x}) \rrbracket .
        \]
    \item More generally, an \vocab{error measure} \(E(g, f)\).
    \item Naturally considered:
    \begin{itemize}
        \item out-of-sample: averaged over unknown \(\bm{x}\);
        \item pointwise: evaluated on one \(\bm{x}\);
        \item classification: \(\llbracket \text{prediction} \neq
            \text{target} \rrbracket\).
    \end{itemize}
\end{itemize}

Here \(\mathcal{E}\) denotes expectation. The three adjectives are three
independent choices that \(\Eout\) happened to make at once, and each could be
made differently.

\begin{moral}
    Classification error \(\llbracket \dots \rrbracket\): often also called
    `0/1 error'.
\end{moral}

\subsection{Pointwise Error Measure}

We can often express \(E(g, f)\) as an average of \(\err(g(\bm{x}),
f(\bm{x}))\), like
\[
    \Eout(g) = \mathop{\mathcal{E}}_{\bm{x} \sim P}
    \underbrace{\llbracket g(\bm{x}) \neq f(\bm{x}) \rrbracket}_{\err(g(\bm{x}),
    f(\bm{x}))} .
\]

\begin{definition}[Pointwise error measure]\label{def:pointwise}
    The function \(\err\) above is called a \vocab{pointwise error measure}.
    With it,
    \[
        \underbrace{\Ein(g) = \frac{1}{N} \sum_{n=1}^N \err(g(\bm{x}_n),
        f(\bm{x}_n))}_{\text{in-sample}}
        \qquad
        \underbrace{\Eout(g) = \mathop{\mathcal{E}}_{\bm{x} \sim P}
        \err(g(\bm{x}), f(\bm{x}))}_{\text{out-of-sample}} .
    \]
\end{definition}

Not every sensible measure is pointwise --- the area under an ROC curve, say,
compares predictions on different inputs --- but we will mainly consider
pointwise \(\err\) for simplicity.

\subsection{Two Important Pointwise Error Measures}

Write the prediction and the target as
\[
    \err\bigl( \underbrace{g(\bm{x})}_{\tilde{y}},
    \underbrace{f(\bm{x})}_{y} \bigr) .
\]

\begin{center}
    \solidrules
    \begin{tabular}{@{}l|l@{}}
        \toprule
        \textbf{0/1 error} & \textbf{squared error}\\
        \midrule
        \(\err(\tilde{y}, y) = \llbracket \tilde{y} \neq y \rrbracket\)
        & \(\err(\tilde{y}, y) = (\tilde{y} - y)^2\)\\
        correct or incorrect? & how far is \(\tilde{y}\) from \(y\)?\\
        often for classification & often for regression\\
        \bottomrule
    \end{tabular}
\end{center}

\begin{question}
    How does \(\err\) `guide' learning?
\end{question}

\subsection{Ideal Mini-Target}

The interplay between noise and error:

\begin{moral}
    \(P(y \mid \bm{x})\) and \(\err\) define the ideal mini-target
    \(f(\bm{x})\).
\end{moral}

With noise there is no single `true' label to recover at \(\bm{x}\): what
nature gives is the whole distribution \(P(y \mid \bm{x})\). The \(f\) we aim
at is therefore not a target function handed to us but the best
\emph{prediction} against that distribution, and `best' only means something
once \(\err\) is fixed.

\begin{definition}[Ideal mini-target]\label{def:mini-target}
    Given \(P(y \mid \bm{x})\) and \(\err\), the \vocab{ideal mini-target}
    at \(\bm{x}\) is the prediction of smallest expected error,
    \[
        f(\bm{x}) = \argmin_{\tilde{y}} \;
        \mathop{\mathcal{E}}_{y \sim P(y \mid \bm{x})} \err(\tilde{y}, y) .
    \]
\end{definition}

So from here on \(f\) denotes the error-minimising function, not a function
that generated the labels; for a deterministic target (\(P(y \mid \bm{x}) =
1\) at \(y = f(\bm{x})\)) and any \(\err\) with \(\err(y, y) = 0 \le
\err(\tilde{y}, y)\), the two coincide.

Take one input whose labels are
\[
    P(y = 1 \mid \bm{x}) = 0.2, \quad P(y = 2 \mid \bm{x}) = 0.7, \quad
    P(y = 3 \mid \bm{x}) = 0.1 ,
\]
and ask which prediction \(\tilde{y}\) has the smallest average error under
each measure:

\begin{center}
    \solidrules
    \begin{tabular}{@{}c|l|l@{}}
        \toprule
        \(\tilde{y}\) & \(\err(\tilde{y}, y) = \llbracket \tilde{y} \neq y
        \rrbracket\) & \(\err(\tilde{y}, y) = (\tilde{y} - y)^2\)\\
        \midrule
        \(1\) & avg.\ \(\err\) \(0.8\) & avg.\ \(\err\) \(1.1\)\\
        \(2\) & avg.\ \(\err\) \(0.3\) \((*)\) & avg.\ \(\err\) \(0.3\)\\
        \(3\) & avg.\ \(\err\) \(0.9\) & avg.\ \(\err\) \(1.5\)\\
        \(1.9\) & avg.\ \(\err\) \(1.0\) (\textcolor{orange!85!black}{really?
        :-)}) & avg.\ \(\err\) \(0.29\) \((*)\)\\
        \midrule
        & \cellcolor{vcregion}\(f(\bm{x}) = \displaystyle\argmax_{y \in
        \mathcal{Y}} P(y \mid \bm{x})\)
        & \cellcolor{vcregion}\(f(\bm{x}) = \displaystyle\sum_{y \in
        \mathcal{Y}} y \cdot P(y \mid \bm{x})\)\\
        \bottomrule
    \end{tabular}
\end{center}

The two starred rows are no accident: they are the two minimisers below.

\begin{proposition}[Ideal mini-targets of 0/1 and squared error]
    \label{prop:mini-target}
    Fix \(\bm{x}\) and write \(\mathcal{E}\) for expectation over \(y \sim
    P(y \mid \bm{x})\).
    \begin{enumerate}
        \item For \(\err(\tilde{y}, y) = \llbracket \tilde{y} \neq y
            \rrbracket\), \(f(\bm{x}) = \argmax_{y \in \mathcal{Y}} P(y \mid
            \bm{x})\) (any maximiser, if there are ties).
        \item For \(\err(\tilde{y}, y) = (\tilde{y} - y)^2\), \(f(\bm{x}) =
            \sum_{y \in \mathcal{Y}} y \cdot P(y \mid \bm{x}) = \mathcal{E}[y]\),
            and this minimiser is unique.
    \end{enumerate}
\end{proposition}

\begin{proof}
    (1) For any real \(\tilde{y}\),
    \[
        \mathcal{E} \llbracket \tilde{y} \neq y \rrbracket
        = \mathbb{P}[y \neq \tilde{y}]
        = 1 - P(y = \tilde{y} \mid \bm{x}) ,
    \]
    where \(P(y = \tilde{y} \mid \bm{x}) = 0\) if \(\tilde{y} \notin
    \mathcal{Y}\). Minimising the left side is maximising \(P(y = \tilde{y}
    \mid \bm{x})\). A \(\tilde{y}\) outside \(\mathcal{Y}\) scores the worst
    possible value \(1\), so the minimum is attained exactly at the labels of
    largest probability.

    (2) Let \(\mu = \mathcal{E}[y]\). Add and subtract \(\mu\):
    \[
        \mathcal{E} (\tilde{y} - y)^2
        = \mathcal{E} \bigl[ (\tilde{y} - \mu) + (\mu - y) \bigr]^2
        = (\tilde{y} - \mu)^2 + 2 (\tilde{y} - \mu)\, \mathcal{E}[\mu - y]
        + \mathcal{E} (\mu - y)^2 .
    \]
    The middle term vanishes because \(\mathcal{E}[\mu - y] = \mu - \mu = 0\),
    and the last term is \(\Var(y)\), which does not depend on \(\tilde{y}\).
    Hence
    \[
        \mathcal{E} (\tilde{y} - y)^2 = (\tilde{y} - \mu)^2 + \Var(y)
        \ge \Var(y) ,
    \]
    with equality if and only if \(\tilde{y} = \mu\).
\end{proof}

On the example, (1) gives \(1 - 0.7 = 0.3\) at \(\tilde{y} = 2\), and \(1\) at
\(\tilde{y} = 1.9\), which is never a label and so is always wrong. (2) gives
\(\mu = 0.2 + 1.4 + 0.3 = 1.9\) and \(\Var(y) = 3.9 - 1.9^2 = 0.29\), the
table's minimum; predicting \(2\) instead costs \((2 - 1.9)^2 + 0.29 = 0.3\).
Same noise, different \(\err\), different ideal mini-target.

\subsection{Learning Flow with Error Measure}

The error measure therefore joins the flow as a component: it is what `\(g
\approx f\)' is judged by.

\begin{center}
    \begin{adjustbox}{max width=\textwidth}
        \begin{tikzpicture}
            \node[mlbox, minimum width=4.4cm, draw=orange!85!black,
                fill=orange!8] (f) at (0,2.5)
                {unknown target\\
                 distribution \(P(y \mid \bm{x})\)\\
                 containing \(f(\bm{x})\) \(+\) noise\\[2pt]
                 {\scriptsize (ideal credit approval formula)}};
            \node[mlbox, minimum width=2.6cm] (P) at (5.4,2.5)
                {unknown\\\(P\) on \(\mathcal{X}\)};
            \node[mlbox, minimum width=5.2cm] (D) at (0,0)
                {training examples\\
                 \(\mathcal{D}\colon (\bm{x}_1, y_1), \dots , (\bm{x}_N, y_N)\)\\[2pt]
                 {\scriptsize (historical records in bank)}};
            \node[mlop, minimum width=2.4cm] (A) at (5.4,0)
                {learning\\algorithm\\\(\mathcal{A}\)};
            \node[mlbox, minimum width=4.4cm] (g) at (10.2,0)
                {final hypothesis\\\(g \approx f\)\\[2pt]
                 {\scriptsize (`learned' formula to be used)}};
            \node[mlbox, minimum width=3.4cm] (H) at (5.4,-2.4)
                {hypothesis set\\\(\mathcal{H}\)\\[2pt]
                 {\scriptsize (set of candidate formulas)}};
            \node[mlbox, minimum width=2.6cm, draw=blue!80!black, very thick]
                (e) at (10.2,-2.4) {error measure\\\(\err\)};
            \draw[mlarrow, dashed, draw=orange!85!black] (f) -- (D)
                node[midway, left, font=\scriptsize, text=orange!85!black]
                {\(y_1, y_2, \dots , y_N\)};
            \draw[mlarrow, dashed, draw=orange!85!black] (f.north)
                .. controls (1,4.4) and (11.1,4.4) .. ([xshift=0.9cm]g.north)
                node[pos=0.5, above, font=\scriptsize, text=orange!85!black]
                {\(y\)};
            \draw[mlarrow, dashed] (P) -- (D)
                node[pos=0.58, font=\scriptsize, fill=white, inner sep=1.5pt]
                {\(\bm{x}_1, \bm{x}_2, \dots , \bm{x}_N\)};
            \draw[mlarrow, dashed] (P.east) -| ([xshift=-0.9cm]g.north)
                node[pos=0.3, above, font=\scriptsize] {\(\bm{x}\)};
            \draw[mlarrow] (D) -- (A);
            \draw[mlarrow] (A) -- (g);
            \draw[mlarrow] (H) -- (A);
            \draw[mlarrow, draw=blue!80!black] (e) -- (g);
        \end{tikzpicture}
    \end{adjustbox}
\end{center}

\begin{moral}
    Extended VC theory/`philosophy' works for most \(\mathcal{H}\) and
    \(\err\).
\end{moral}

The proof in \autoref{sec:pictorial} was for 0/1 error; analogous bounds exist
for other bounded error measures, with \(\dvc\) replaced by suitable
generalisations. We take the philosophy on trust from here.

\begin{exercise}[Fun Time]
    Consider the following \(P(y \mid \bm{x})\) and \(\err(\tilde{y}, y) =
    \abs{\tilde{y} - y}\). Which of the following is the ideal mini-target
    \(f(\bm{x})\)?
    \[
        \begin{gathered}
            P(y = 1 \mid \bm{x}) = 0.10, \quad P(y = 2 \mid \bm{x}) = 0.35,\\
            P(y = 3 \mid \bm{x}) = 0.15, \quad P(y = 4 \mid \bm{x}) = 0.40.
        \end{gathered}
    \]
    \begin{enumerate}[(1)]
        \item \(2.5 =\) average within \(\mathcal{Y} = \left\{ 1, 2, 3, 4
            \right\}\);
        \item \(2.85 =\) weighted mean from \(P(y \mid \bm{x})\);
        \item \(3 =\) weighted median from \(P(y \mid \bm{x})\);
        \item \(4 = \argmax P(y \mid \bm{x})\).
    \end{enumerate}
\end{exercise}

\begin{proof}[Reference Answer]
    (3). For the `absolute error', the weighted median provably results in the
    minimum average \(\err\). The average errors of the four candidates are
    \(1.0\), \(0.965\), \(0.95\) and \(1.15\). Moving \(\tilde{y}\) to the
    right by a little changes the average by (mass to the left) \(-\) (mass to
    the right); that is negative until half the mass is on the left, which
    happens at \(3\) (\(0.10 + 0.35 + 0.15 = 0.6 \ge 0.5\)).
\end{proof}

\section{Algorithmic Error Measure}\label{sec:algo-error}

\source{slides p.14--19}

The last section treated \(\err\) as given. Who gives it?

\subsection{Choice of Error Measure}

\paragraph{Fingerprint verification.} A fingerprint goes in, and the target says
whether it belongs to you:

\begin{center}
    \begin{tikzpicture}
        \node (fp) at (0,0)
            {\includegraphics[width=1.3cm]{Figures/fingerprint.jpg}};
        \node[mlbox, minimum width=1.6cm, minimum height=7mm] (f) at (2.6,0)
            {\(f\)};
        \draw[mlarrow] (fp) -- (f);
        \draw[mlarrow] (f) -- (4.2,0);
        \node[anchor=west] at (4.2,0)
            {\(\left\{ \begin{array}{@{}ll@{}}
                +1 & \textcolor{ForestGreen}{\text{you}}\\
                -1 & \textcolor{red!80!black}{\text{intruder}}
            \end{array} \right.\)};
    \end{tikzpicture}
\end{center}

There are two types of error: \textcolor{orange!85!black}{false accept} and
\textcolor{Purple}{false reject}.

\begin{center}
    \solidrules
    \begin{tabular}{@{}cc|cc@{}}
        \toprule
        & & \multicolumn{2}{c}{\(g\)}\\
        & & \(+1\) & \(-1\)\\
        \midrule
        \multirow{2}{*}{\(f\)} & \(+1\) & \textcolor{black!45}{no error}
        & \textcolor{Purple}{false reject}\\
        & \(-1\) & \textcolor{orange!85!black}{false accept}
        & \textcolor{black!45}{no error}\\
        \bottomrule
    \end{tabular}
\end{center}

\begin{moral}
    0/1 error penalizes both types equally.
\end{moral}

Whether that is right depends entirely on who uses the verifier.

\subsection{Fingerprint Verification for Supermarket}

\begin{center}
    \solidrules
    \begin{tabular}{@{}cc|cc@{}}
        \toprule
        & & \multicolumn{2}{c}{\(g\)}\\
        & & \(+1\) & \(-1\)\\
        \midrule
        \multirow{2}{*}{\(f\)} & \(+1\) & \textcolor{black!45}{\(0\)}
        & \textcolor{Purple}{\(10\)}\\
        & \(-1\) & \textcolor{orange!85!black}{\(1\)}
        & \textcolor{black!45}{\(0\)}\\
        \bottomrule
    \end{tabular}
\end{center}

\begin{itemize}
    \item Supermarket: fingerprint for discount.
    \item False reject: \textbf{very unhappy customer, lose future business}.
    \item False accept: give away a minor discount, intruder left fingerprint
        :-)
\end{itemize}

\subsection{Fingerprint Verification for CIA}

\begin{center}
    \solidrules
    \begin{tabular}{@{}cc|cc@{}}
        \toprule
        & & \multicolumn{2}{c}{\(g\)}\\
        & & \(+1\) & \(-1\)\\
        \midrule
        \multirow{2}{*}{\(f\)} & \(+1\) & \textcolor{black!45}{\(0\)}
        & \textcolor{Purple}{\(1\)}\\
        & \(-1\) & \textcolor{orange!85!black}{\(1000\)}
        & \textcolor{black!45}{\(0\)}\\
        \bottomrule
    \end{tabular}
\end{center}

\begin{itemize}
    \item CIA: fingerprint for entrance.
    \item False accept: \textbf{very serious consequences!}
    \item False reject: unhappy employee, but so what? :-)
\end{itemize}

Same inputs, same target, same two kinds of mistake --- and cost ratios of
\(10 : 1\) one way and \(1 : 1000\) the other. The best \(g\) for one
application is a poor \(g\) for the other.

\subsection{Take-home Message for Now}

\begin{moral}
    \(\err\) is application/user-dependent.
\end{moral}

In practice the user rarely writes the true \(\err\) down, so the algorithm
works with a stand-in.

\begin{definition}[Algorithmic error measure]\label{def:errhat}
    An \vocab{algorithmic error measure} \(\errhat\) is the error measure
    \(\mathcal{A}\) actually uses in place of \(\err\). It may be
    \begin{itemize}
        \item true: just \(\err\);
        \item plausible:
        \begin{itemize}
            \item 0/1: minimum `flipping noise' --- NP-hard to optimize,
                remember? :-)
            \item squared: minimum Gaussian noise;
        \end{itemize}
        \item friendly: easy to optimize for \(\mathcal{A}\):
        \begin{itemize}
            \item closed-form solution;
            \item convex objective function.
        \end{itemize}
    \end{itemize}
\end{definition}

`Plausible' means justified by a belief about the noise: if labels are flipped
at random, the best one can do is make the fewest disagreements, which is 0/1
error --- and minimising it over perceptrons is the NP-hard problem of
\autoref{sec:nonsep}; if the noise is Gaussian, maximum likelihood is least
squares. `Friendly' means justified by what the optimiser can do.

\begin{moral}
    \(\errhat\): more in next lectures.
\end{moral}

\subsection{Learning Flow with Algorithmic Error Measure}

\begin{center}
    \begin{adjustbox}{max width=\textwidth}
        \begin{tikzpicture}
            \node[mlbox, minimum width=4.4cm, draw=orange!85!black,
                fill=orange!8] (f) at (0,2.5)
                {unknown target\\
                 distribution \(P(y \mid \bm{x})\)\\
                 containing \(f(\bm{x})\) \(+\) noise\\[2pt]
                 {\scriptsize (ideal credit approval formula)}};
            \node[mlbox, minimum width=2.6cm] (P) at (5.4,2.5)
                {unknown\\\(P\) on \(\mathcal{X}\)};
            \node[mlbox, minimum width=5.2cm] (D) at (0,0)
                {training examples\\
                 \(\mathcal{D}\colon (\bm{x}_1, y_1), \dots , (\bm{x}_N, y_N)\)\\[2pt]
                 {\scriptsize (historical records in bank)}};
            \node[mlop, minimum width=2.4cm] (A) at (5.4,0)
                {learning\\algorithm\\\(\mathcal{A}\)\\
                 \colorbox{Purple!18}{\(\errhat\)}};
            \node[mlbox, minimum width=4.4cm] (g) at (10.2,0)
                {final hypothesis\\\(g \approx f\)\\[2pt]
                 {\scriptsize (`learned' formula to be used)}};
            \node[mlbox, minimum width=3.4cm] (H) at (5.4,-2.6)
                {hypothesis set\\\(\mathcal{H}\)\\[2pt]
                 {\scriptsize (set of candidate formulas)}};
            \node[mlbox, minimum width=2.6cm] (e) at (10.2,-2.6)
                {error measure\\\(\err\)};
            \draw[mlarrow, dashed, draw=orange!85!black] (f) -- (D)
                node[midway, left, font=\scriptsize, text=orange!85!black]
                {\(y_1, y_2, \dots , y_N\)};
            \draw[mlarrow, dashed, draw=orange!85!black] (f.north)
                .. controls (1,4.4) and (11.1,4.4) .. ([xshift=0.9cm]g.north)
                node[pos=0.5, above, font=\scriptsize, text=orange!85!black]
                {\(y\)};
            \draw[mlarrow, dashed] (P) -- (D)
                node[pos=0.58, font=\scriptsize, fill=white, inner sep=1.5pt]
                {\(\bm{x}_1, \bm{x}_2, \dots , \bm{x}_N\)};
            \draw[mlarrow, dashed] (P.east) -| ([xshift=-0.9cm]g.north)
                node[pos=0.3, above, font=\scriptsize] {\(\bm{x}\)};
            \draw[mlarrow] (D) -- (A);
            \draw[mlarrow] (A) -- (g);
            \draw[mlarrow] (H) -- (A);
            \draw[mlarrow, draw=blue!80!black] (e) -- (g);
        \end{tikzpicture}
    \end{adjustbox}
\end{center}

\begin{moral}
    \(\err\): application goal; \(\errhat\): a key part of many
    \(\mathcal{A}\).
\end{moral}

\(\err\) sits outside the algorithm and judges \(g\); \(\errhat\) sits inside and
steers the search. The art is choosing \(\errhat\) close enough to \(\err\) to
aim at the right goal, and friendly enough to be optimised.

\begin{exercise}[Fun Time]
    Consider \(\err\) below for CIA. What is \(\Ein(g)\) when using this
    \(\err\)?
    \begin{center}
        \solidrules
        \begin{tabular}{@{}cc|cc@{}}
            \toprule
            & & \multicolumn{2}{c}{\(g\)}\\
            & & \(+1\) & \(-1\)\\
            \midrule
            \multirow{2}{*}{\(f\)} & \(+1\) & \(0\) & \(1\)\\
            & \(-1\) & \(1000\) & \(0\)\\
            \bottomrule
        \end{tabular}
    \end{center}
    \begin{enumerate}[(1)]
        \item \(\displaystyle \frac{1}{N} \sum_{n=1}^N \llbracket y_n \neq
            g(\bm{x}_n) \rrbracket\);
        \item \(\displaystyle \frac{1}{N} \left( \sum_{y_n = +1} \llbracket
            y_n \neq g(\bm{x}_n) \rrbracket + 1000 \sum_{y_n = -1} \llbracket
            y_n \neq g(\bm{x}_n) \rrbracket \right)\);
        \item \(\displaystyle \frac{1}{N} \left( \sum_{y_n = +1} \llbracket
            y_n \neq g(\bm{x}_n) \rrbracket - 1000 \sum_{y_n = -1} \llbracket
            y_n \neq g(\bm{x}_n) \rrbracket \right)\);
        \item \(\displaystyle \frac{1}{N} \left( 1000 \sum_{y_n = +1}
            \llbracket y_n \neq g(\bm{x}_n) \rrbracket + \sum_{y_n = -1}
            \llbracket y_n \neq g(\bm{x}_n) \rrbracket \right)\).
    \end{enumerate}
\end{exercise}

\begin{proof}[Reference Answer]
    (2). When \(y_n = -1\), the false positive made on such \((\bm{x}_n,
    y_n)\) is penalized \(1000\) times more!
\end{proof}

\section{Weighted Classification}\label{sec:weighted}

\source{slides p.20--24}

The CIA error of the quiz is a concrete, useful \(\err\) that is not 0/1. This
section shows how to learn with it using nothing but the tools of
\autoref{sec:nonsep}.

\subsection{Weighted Classification}

\paragraph{CIA cost (error, loss, \dots) matrix.}
\begin{center}
    \solidrules
    \begin{tabular}{@{}cc|cc@{}}
        \toprule
        & & \multicolumn{2}{c}{\(h(\bm{x})\)}\\
        & & \(+1\) & \(-1\)\\
        \midrule
        \multirow{2}{*}{\(y\)} & \(+1\) & \(0\) & \(1\)\\
        & \(-1\) & \(1000\) & \(0\)\\
        \bottomrule
    \end{tabular}
\end{center}

Out-of-sample,
\[
    \Eout(h) = \mathop{\mathcal{E}}_{(\bm{x}, y) \sim P}
    \left\{ \begin{array}{@{}ll@{}}
        1 & \text{if } y = +1\\
        1000 & \text{if } y = -1
    \end{array} \right\} \cdot \llbracket y \neq h(\bm{x}) \rrbracket ;
\]
in-sample,
\[
    \Ein(h) = \frac{1}{N} \sum_{n=1}^N
    \left\{ \begin{array}{@{}ll@{}}
        1 & \text{if } y_n = +1\\
        1000 & \text{if } y_n = -1
    \end{array} \right\} \cdot \llbracket y_n \neq h(\bm{x}_n) \rrbracket .
\]
The cost of a mistake depends only on the true label, so the matrix collapses
into one weight per example.

\begin{definition}[Weighted classification]\label{def:weighted}
    \vocab{Weighted classification} is classification with a different
    `weight' for different \((\bm{x}, y)\): the pointwise error is \(w(\bm{x},
    y) \cdot \llbracket y \neq h(\bm{x}) \rrbracket\) for a given weight
    \(w(\bm{x}, y) \ge 0\).
\end{definition}

\subsection{Minimizing \texorpdfstring{\(\Ein\)}{Ein} for Weighted
Classification}

Write \(\Ein^{\mathrm{w}}\) for the weighted in-sample error,
\[
    \Ein^{\mathrm{w}}(h) = \frac{1}{N} \sum_{n=1}^N
    \left\{ \begin{array}{@{}ll@{}}
        1 & \text{if } y_n = +1\\
        1000 & \text{if } y_n = -1
    \end{array} \right\} \cdot \llbracket y_n \neq h(\bm{x}_n) \rrbracket ,
\]
and \(\Ein^{0/1}\) for the plain fraction of mistakes.

\paragraph{Naïve thoughts.}
\begin{itemize}
    \item PLA: doesn't matter if linear separable. :-) Then some \(\bm{w}\)
        has \(\Ein^{\mathrm{w}} = 0\), whatever the weights, and PLA finds it.
    \item Pocket: modify the pocket-replacement rule --- if \(\bm{w}_{t+1}\)
        reaches smaller \(\Ein^{\mathrm{w}}\) than \(\hat{\bm{w}}\), replace
        \(\hat{\bm{w}}\) by \(\bm{w}_{t+1}\).
\end{itemize}

\begin{question}
    Pocket has some guarantee on \(\Ein^{0/1}\); does the modified pocket have
    a similar guarantee on \(\Ein^{\mathrm{w}}\)?
\end{question}

\subsection{Systematic Route: Connect
\texorpdfstring{\(\Ein^{\mathrm{w}}\)}{Einw} and
\texorpdfstring{\(\Ein^{0/1}\)}{Ein01}}

Rather than re-prove anything about pocket, turn the weighted problem into an
unweighted one.

\begin{center}
    \begin{minipage}[t]{0.44\textwidth}
        \centering
        \textbf{original problem}\\[3pt]
        \solidrules
        \begin{tabular}{@{}cc|cc@{}}
            \toprule
            & & \multicolumn{2}{c}{\(h(\bm{x})\)}\\
            & & \(+1\) & \(-1\)\\
            \midrule
            \multirow{2}{*}{\(y\)} & \(+1\) & \(0\) & \(1\)\\
            & \(-1\) & \textcolor{orange!85!black}{\(1000\)} & \(0\)\\
            \bottomrule
        \end{tabular}\\[6pt]
        \(\mathcal{D}\colon\)
        \(\begin{array}{c}
            (\bm{x}_1, +1)\\
            \textcolor{red!80!black}{(\bm{x}_2, -1)}\\
            \textcolor{red!80!black}{(\bm{x}_3, -1)}\\
            \dots\\
            (\bm{x}_{N-1}, +1)\\
            (\bm{x}_N, +1)
        \end{array}\)
    \end{minipage}
    \begin{minipage}[t]{0.54\textwidth}
        \centering
        \textbf{equivalent problem}\\[3pt]
        \solidrules
        \begin{tabular}{@{}cc|cc@{}}
            \toprule
            & & \multicolumn{2}{c}{\(h(\bm{x})\)}\\
            & & \(+1\) & \(-1\)\\
            \midrule
            \multirow{2}{*}{\(y\)} & \(+1\) & \(0\) & \(1\)\\
            & \(-1\) & \(1\) & \(0\)\\
            \bottomrule
        \end{tabular}\\[6pt]
        \(\begin{array}{c}
            (\bm{x}_1, +1)\\
            \textcolor{red!80!black}{(\bm{x}_2, -1), (\bm{x}_2, -1), \dots ,
            (\bm{x}_2, -1)}\\
            \textcolor{red!80!black}{(\bm{x}_3, -1), (\bm{x}_3, -1), \dots ,
            (\bm{x}_3, -1)}\\
            \dots\\
            (\bm{x}_{N-1}, +1)\\
            (\bm{x}_N, +1)
        \end{array}\)
    \end{minipage}
\end{center}

\begin{moral}
    After copying \(-1\) examples \(1000\) times, \(\Ein^{\mathrm{w}}\) for LHS
    \(\equiv\) \(\Ein^{0/1}\) for RHS!
\end{moral}

A mistake on \((\bm{x}_2, -1)\) costs \(1000\) on the left and is repeated on
all \(1000\) copies on the right, so the two error counts agree term by term
for every \(h\). (The normalisations differ, \(N\) against the size of the
copied set, but that is a constant factor and does not change which \(h\)
minimises.) Pocket's guarantee for \(\Ein^{0/1}\) on the copied data therefore
\emph{is} a guarantee for \(\Ein^{\mathrm{w}}\) on the original.

\subsection{Weighted Pocket Algorithm}

Copying \(1000\) times is wasteful; the copies only need to exist in how often
examples are picked and how errors are counted.

\begin{center}
    \solidrules
    \begin{tikzpicture}[scale=0.8]
        \fill[poshalf] (0,3.2) -- (3.2,0) -- (3.2,3.2) -- cycle;
        \fill[neghalf] (0,3.2) -- (3.2,0) -- (0,0) -- cycle;
        \draw[sepline] (0,3.2) -- (3.2,0);
        \draw[black!45] (0,0) rectangle (3.2,3.2);
        \foreach \p in {(1.5,2.4), (2.2,2.8), (2.6,1.6), (1.9,1.9),
            (2.9,2.4), (0.4,3.0)}{\node[dpos] at \p {};}
        \foreach \p in {(0.5,0.9), (1.2,1.0), (0.8,1.9), (1.8,0.6),
            (0.4,2.0), (2.7,0.3)}{\node[dneg] at \p {\(\times\)};}
        \node[dneg] at (2.4,2.1) {\(\times\)};
        \node[dneg] at (1.3,2.9) {\(\times\)};
        \node[dpos] at (0.9,0.7) {};
        \node[dpos] at (0.3,0.25) {};
    \end{tikzpicture}
    \qquad
    \begin{tabular}[b]{@{}cc|cc@{}}
        \toprule
        & & \multicolumn{2}{c}{\(h(\bm{x})\)}\\
        & & \(+1\) & \(-1\)\\
        \midrule
        \multirow{2}{*}{\(y\)} & \(+1\) & \(0\) & \(1\)\\
        & \(-1\) & \(1000\) & \(0\)\\
        \bottomrule
    \end{tabular}
\end{center}

Using `virtual copying', the \vocab{weighted pocket algorithm} includes:

\begin{itemize}
    \item weighted PLA: randomly check \(-1\) example mistakes with \(1000\)
        times more probability;
    \item weighted pocket replacement: if \(\bm{w}_{t+1}\) reaches smaller
        \(\Ein^{\mathrm{w}}\) than \(\hat{\bm{w}}\), replace \(\hat{\bm{w}}\) by
        \(\bm{w}_{t+1}\).
\end{itemize}

Against \autoref{alg:pocket}, only two lines change:

\begin{algorithm}[H]
    \caption{Weighted Pocket Algorithm}\label{alg:weighted-pocket}
    \KwIn{training examples \(\mathcal{D}\), weights \(1\) for \(y_n = +1\) and
        \(1000\) for \(y_n = -1\), an iteration budget}
    \KwOut{\(\hat{\bm{w}}\), returned as \(g\)}
    initialize the pocket weights \(\hat{\bm{w}}\)\;
    \For{\(t = 0, 1, \dots\) until enough iterations}{
        find a random mistake of \(\bm{w}_t\), \((\bm{x}_{n(t)}, y_{n(t)})\),
        picking each mistake with probability proportional to its weight\;
        (try to) correct the mistake by
        \(\bm{w}_{t+1} \gets \bm{w}_t + y_{n(t)} \bm{x}_{n(t)}\)\;
        \lIf{\(\Ein^{\mathrm{w}}(\bm{w}_{t+1}) < \Ein^{\mathrm{w}}(\hat{\bm{w}})\)}
            {\(\hat{\bm{w}} \gets \bm{w}_{t+1}\)}
    }
    \Return \(\hat{\bm{w}}\) as \(g\)\;
\end{algorithm}

Picking a random mistake of the copied data set is exactly picking an original
mistake with probability proportional to its weight, and the copied data's
\(\Ein^{0/1}\) is \(\Ein^{\mathrm{w}}\) up to a constant; so this \emph{is}
pocket on the copied data, run without making the copies.

\begin{moral}
    Systematic route (called `reduction'): can be applied to many other
    algorithms!
\end{moral}

\begin{exercise}[Fun Time]
    Consider the CIA cost matrix. If there are \(10\) examples with
    \(y_n = -1\) (intruder) and \(999{,}990\) examples with \(y_n = +1\)
    (you). What would \(\Ein^{\mathrm{w}}(h)\) be for a constant \(h(\bm{x})\)
    that always returns \(+1\)?
    \begin{center}
        \solidrules
        \begin{tabular}{@{}cc|cc@{}}
            \toprule
            & & \multicolumn{2}{c}{\(h(\bm{x})\)}\\
            & & \(+1\) & \(-1\)\\
            \midrule
            \multirow{2}{*}{\(y\)} & \(+1\) & \(0\) & \(1\)\\
            & \(-1\) & \(1000\) & \(0\)\\
            \bottomrule
        \end{tabular}
        \qquad
        \begin{enumerate*}[(1)]
            \item \(0.001\); \item \(0.01\); \item \(0.1\); \item \(1\).
        \end{enumerate*}
    \end{center}
\end{exercise}

\begin{proof}[Reference Answer]
    (2). The constant \(h\) is wrong on exactly the \(10\) intruders, each
    costing \(1000\), out of \(N = 10^6\) examples: \(\Ein^{\mathrm{w}} = 10
    \cdot 1000 / 10^6 = 0.01\). While the quiz is a simple evaluation, it is not
    uncommon that the data is very unbalanced for such an application. Properly
    `setting' the weights can be used to avoid the lazy constant prediction:
    under plain 0/1 error the same \(h\) would score \(10^{-5}\), and look
    nearly perfect.
\end{proof}

\hrulebar

\subsection*{Summary}

\begin{description}
    \item[Noise and Probabilistic Target] can replace \(f(\bm{x})\) by \(P(y
        \mid \bm{x})\).
    \item[Error Measure] affects the `ideal' target.
    \item[Algorithmic Error Measure] user-dependent \(\implies\) plausible or
        friendly.
    \item[Weighted Classification] easily done by virtual `example copying'.
\end{description}

The \emph{Why can machines learn?} part closes here: with noise and a general
error measure, the VC philosophy still says a finite \(\dvc\) and enough data
make \(\Ein\) a trustworthy proxy for \(\Eout\). What remains is to make
\(\Ein\) small, and so far the only tools are PLA and pocket. Next: more
algorithms, please? :-)
